% Standalone problem document, generated from the adjacent README.md.
% Compile with XeLaTeX. Mathematical content is maintained in Markdown.
\documentclass[11pt,a4paper]{article}
\usepackage[fontsize=10.5pt]{scrextend}
\usepackage[margin=22mm,headheight=15pt,headsep=9mm,footskip=11mm]{geometry}
\usepackage{fontspec}
\setmainfont{texgyrepagella}[Extension=.otf,UprightFont=*-regular,BoldFont=*-bold,ItalicFont=*-italic,BoldItalicFont=*-bolditalic]
\setsansfont{texgyreheros}[Extension=.otf,UprightFont=*-regular,BoldFont=*-bold,ItalicFont=*-italic,BoldItalicFont=*-bolditalic]
\setmonofont{texgyrecursor}[Extension=.otf,UprightFont=*-regular,BoldFont=*-bold,ItalicFont=*-italic,BoldItalicFont=*-bolditalic,Scale=MatchLowercase]
\usepackage{amsmath,amssymb}
\usepackage{unicode-math}
\setmathfont{texgyrepagella-math.otf}
\usepackage{xcolor,microtype,enumitem,needspace,fancyhdr}
\usepackage{bookmark}
\definecolor{ink}{HTML}{17344B}
\definecolor{accent}{HTML}{197D80}
\definecolor{muted}{HTML}{596773}
\hypersetup{colorlinks=true,linkcolor=accent,urlcolor=accent,pdftitle={MI-12 - Marcus's
inequality for the permanent of block
permanents},pdfauthor={Open Problems in Numerical Linear Algebra}}
\urlstyle{same}
\setlength{\parindent}{0pt}
\setlength{\parskip}{5pt plus 1pt minus 1pt}
\linespread{1.04}
\setlength{\emergencystretch}{3em}
\widowpenalty=10000
\clubpenalty=10000
\displaywidowpenalty=10000
\allowdisplaybreaks[1]
\setlist{nosep,leftmargin=1.5em}
\providecommand{\tightlist}{\setlength{\itemsep}{0pt}\setlength{\parskip}{0pt}}
\providecommand{\pandocbounded}[1]{#1}
\pagestyle{fancy}
\fancyhf{}
\fancyhead[L]{\sffamily\footnotesize\color{muted}OPEN PROBLEMS IN NUMERICAL LINEAR ALGEBRA}
\fancyhead[R]{\sffamily\bfseries\footnotesize\color{accent}MI-12}
\fancyfoot[L]{\parbox[b]{0.9\textwidth}{\sffamily\fontsize{8}{10}\selectfont\color{muted}Editorial ratings; a bounded search cannot certify that a problem remains unsolved.\\Literature check: 2026-09-08}}
\fancyfoot[R]{\sffamily\footnotesize\color{muted}\thepage}
\renewcommand{\headrulewidth}{0.3pt}
\renewcommand{\headrule}{\hbox to\headwidth{\color{accent}\leaders\hrule height \headrulewidth\hfill}}
\makeatletter
\renewcommand{\section}{\Needspace{5\baselineskip}\@startsection{section}{1}{0pt}{12pt plus 2pt minus 2pt}{4pt}{\sffamily\bfseries\large\color{ink}}}
\renewcommand{\subsection}{\Needspace{4\baselineskip}\@startsection{subsection}{2}{0pt}{9pt plus 1pt minus 1pt}{3pt}{\sffamily\bfseries\normalsize\color{ink}}}
\makeatother
\setcounter{secnumdepth}{0}
\begin{document}
{\sffamily\small\color{muted}Matrix inequalities and norms\par}
\vspace{3pt}
{\raggedright\hyphenpenalty=10000\exhyphenpenalty=10000\sffamily\bfseries\fontsize{21}{24}\selectfont\color{ink}Marcus's
inequality for the permanent of block permanents\par}
\vspace{7pt}
{\sffamily\small\textbf{Difficulty:} extreme\quad\textbf{Importance:} interesting
to specialist\par}
\vspace{4pt}
{\color{accent}\hrule height 0.8pt}
\vspace{6pt}
\textbf{Status:} explicit conjectured inequality; no later resolution
located

For a square matrix \(C\) of order \(r\), write \[
\operatorname{per}C=\sum_{\sigma\in S_r}\prod_{i=1}^r c_{i,\sigma(i)}.
\] Let \(m,k\ge2\) be arbitrary integers and let
\(A\in\mathbb C^{mk\times mk}\) be Hermitian positive semidefinite.
Partition \(A\) into an \(m\times m\) array of \(k\times k\) blocks
\(A_{ij}\), and form \[
G=(g_{ij})_{i,j=1}^m,\qquad g_{ij}=\operatorname{per}(A_{ij}).
\] Must \[
\operatorname{per}A\ge\operatorname{per}G
\] always hold?

\subsection{Relevance}\label{relevance}

This asks how a costly matrix polynomial behaves under block aggregation
by the same polynomial. It is a precise constraint on hierarchical
permanent computation for PSD matrices; positivity is assumed for the
whole matrix, not separately for off-diagonal blocks.

\subsection{References}\label{references}

\begin{itemize}
\tightlist
\item
  I. M. Wanless, \emph{Lieb's permanental dominance conjecture} (2022),
  Conjecture 3 and the following paragraph
  (\href{https://arxiv.org/pdf/2202.01867}{primary manuscript}).
\item
  F. Zhang, \emph{An update on a few permanent conjectures}, Special
  Matrices 4 (2016), 305--316, discussion of Marcus's
  permanent-of-permanents conjecture
  (\href{https://arxiv.org/pdf/1608.02844}{primary manuscript};
  \href{https://doi.org/10.1515/spma-2016-0030}{journal}).
\item
  E. H. Lieb, \emph{Proofs of some Conjectures on Permanents}, Journal
  of Mathematics and Mechanics 16 (1966), 127--134, the two-block
  inequality
  (\href{https://iumj.s3-us-west-2.amazonaws.com/abstracts/16008_abs.pdf}{publisher's
  first page}).
\end{itemize}

\subsection{Status check --- 2026-09-08}\label{status-check-2026-09-08}

Wanless records this inequality as open and identifies Lieb's proof for
\(m=2\). Searches for ``Marcus conjecture'', ``permanent of
permanents'', ``block permanents'', ``proof'', and 2025--2026 found no
general resolution. This entry is specifically the inequality component:
some historical formulations additionally conjecture an equality
characterization. The trivial block sizes \(m=1\) and \(k=1\) are
omitted, and no unqualified equality characterization is imported.
Lieb's general permanental dominance conjecture would imply this
assertion, but this is a separately stated historical conjecture.
\end{document}
